\documentclass[11pt,a4paper]{article}

\usepackage[utf8]{inputenc}
\usepackage[T1]{fontenc}
\usepackage{amsmath,amssymb,bm}
\usepackage{booktabs}
\usepackage[margin=2.5cm]{geometry}
\usepackage{hyperref}
\usepackage{enumitem}
\usepackage{listings}
\usepackage{xcolor}

\hypersetup{
    colorlinks=true,
    linkcolor=blue,
    citecolor=blue,
    urlcolor=blue
}

\lstset{
    basicstyle=\ttfamily\footnotesize,
    columns=fullflexible,
    breaklines=true,
    breakatwhitespace=true,
    frame=single,
    framerule=0.3pt,
    rulecolor=\color{gray!50},
    backgroundcolor=\color{gray!5},
    xleftmargin=0.5em,
    xrightmargin=0.5em,
    aboveskip=0.6em,
    belowskip=0.6em
}

\newcommand{\mb}[1]{\mathbf{#1}}
\newcommand{\Heff}{\mb{H}_{\text{eff}}}
\newcommand{\hth}{\mb{h}_{\text{th}}}
\newcommand{\Ms}{M_s}
\newcommand{\Aex}{A_{\text{ex}}}
\newcommand{\Keff}{K_{\text{eff}}}
\newcommand{\Vcell}{V_{\text{cell}}}
\newcommand{\tCo}{t_{\text{Co}}}
\newcommand{\Rth}{R_{\text{th}}}
\newcommand{\jcur}{j}
\newcommand{\kB}{k_B}

\title{Stochastic LLGS Solver for Thermal Skyrmion Dynamics:\\
Theory, Implementation Plan, and Reasoning}
\author{Rui Barreira}
\date{May 2026}

\begin{document}
\maketitle

\tableofcontents
\bigskip\hrule\bigskip

%======================================================================
\section{Motivation and questions to answer}
%======================================================================

The deterministic LLGS simulator in \texttt{src/simulator/} reproduces
the $T = 0$ dynamics of a SAF skyrmion under SOT drive but cannot
address the following physics questions:

\begin{enumerate}[leftmargin=2em]
    \item[\textbf{Q1}] At room temperature, can the same drive current
    sustain the same skyrmion velocity as at $T = 0$, or must the
    drive be reduced?
    \item[\textbf{Q2}] At high drive velocity, is the directional
    control preserved? In a SAF the topological Hall force cancels
    between layers; thermal fluctuations partially break this
    cancellation. How does the skyrmion Hall angle $\theta_H$ and the
    transverse position variance $\sigma_y(t)$ behave as $T$ and
    $\jcur$ grow?
    \item[\textbf{Q3}] Does the skyrmion size affect velocity and
    interactions? Two complementary "pair" scenarios are in scope,
    both measured:
    \begin{itemize}[leftmargin=1.5em, nosep]
        \item[(Q3a)] \emph{Inter-layer pair}: the top-Co + bot-Co
        AF-coupled pair within a \emph{single} SAF skyrmion. Vary
        $D, H_z$ to span $R$; record top-bot lateral offset
        $|\mb{r}_{c,\text{top}} - \mb{r}_{c,\text{bot}}|$,
        compensation residual $Q_{\text{top}} + Q_{\text{bot}}$, and
        $d_{\text{top}}$ vs $d_{\text{bot}}$ under SOT drive as a
        function of $R$.
        \item[(Q3b)] \emph{Inter-skyrmion pair}: two SAF skyrmions on
        one lattice at separation $r_{\text{init}}$. Reconstruct
        $V(r)$ from ensemble-averaged $r(t)$ at $\jcur = 0$.
    \end{itemize}
    The skyrmion radius $R$ is set by the balance of $\Aex$, $D$,
    $\Keff$ and the applied field $H_z$; the domain-wall width
    $\Delta = \sqrt{\Aex/\Keff}$ sets the natural length scale for
    both the inter-layer coupling and the inter-skyrmion potential.
    \item[\textbf{Q4}] Is there an optimal temperature that maximizes
    drift speed while preserving the skyrmion against thermal
    annihilation?
\end{enumerate}

Answering these requires solving the \emph{stochastic} LLGS equation
of Brown~\cite{brown1963} on the SAF lattice, with Joule heating
coupling the drive current density to the local temperature.

%======================================================================
\section{Theory}
%======================================================================

\subsection{Notation and the fixed-temperature assumption}
\label{sec:notation_fixed_T}

The state variable evolved by the stochastic LLGS equation is the
\emph{unit} magnetization vector
%
\begin{equation}
\mb{m}(\mb{r}, t) = \frac{\mb{M}(\mb{r}, t)}{M_s(T)},
\qquad |\mb{m}| = 1,
\end{equation}
%
where $\mb{M}$ is the magnetization (A/m) and $M_s(T)$ is the
\emph{saturation magnetization}, i.e.\ the maximum magnitude that
$|\mb{M}|$ attains when all atomic moments align. $M_s$ is a function
of temperature, vanishing at the Curie point $T_C$ and obeying Bloch's
law at low temperature:
%
\begin{equation}\label{eq:bloch}
M_s(T) = M_s(0) \left[1 - (T/T_C)^{3/2} + \cdots \right].
\end{equation}
%
The micromagnetic / LLGS framework treats $|\mb{M}|$ and the direction
$\mb{m}$ as separate degrees of freedom that live on widely different
timescales:
%
\begin{itemize}[leftmargin=2em, nosep]
    \item \textbf{Transverse} (direction of $\mb{m}$): GHz, governed
        by precession and damping. This is what LLGS solves for.
    \item \textbf{Longitudinal} ($|\mb{M}|$): THz, set by atomic-scale
        exchange and equilibrating to $M_s(T)$ on every LLGS timestep.
\end{itemize}
%
Under the standard micromagnetic assumption that the longitudinal mode
is equilibrated, $|\mb{M}| = M_s(T)$ \emph{at every instant} and the
constraint $|\mb{m}| = 1$ is exact. This requires
%
\begin{enumerate}[leftmargin=2em, nosep]
    \item $T \ll T_C$, so $M_s(T) \approx M_s(0)$ and material
        parameters $(\Aex, K, D)$ are not strongly $T$-dependent;
    \item $T$ slowly varying compared to LLGS timescales, so $M_s$ can
        be treated as a static constant within one trajectory.
\end{enumerate}
%
The thermal field $\hth$ introduced in Sec.~2.2 \emph{rotates} $\mb{m}$
on the unit sphere and does \emph{not} model longitudinal $|\mb{M}|$
fluctuations: those are absorbed into $M_s(T)$ as a static rescaling.
The Heun scheme used in the integrator (Sec.~3) preserves
$|\mb{m}|=1$ to machine precision at every step.
%
\paragraph{Joule heating.} When Joule heating is enabled
(\texttt{joule\_heating.py}), the lattice temperature $T(t)$ rises over
the course of a trajectory. The implementation re-evaluates the
noise variance $\sigma^2(T)$ at each step (Sec.~2.1) but keeps $M_s$
fixed at its operating-temperature value; updating $M_s(T)$ on the fly
would couple the longitudinal and transverse modes and break the
$|\mb{m}|=1$ constraint, which is outside the scope of LLGS. The
implicit assumption is that the operating-temperature range explored
in a single run is small enough that $M_s(T)$ variation is negligible
relative to the noise-induced dynamics.

\subsection{Brown's stochastic LLG equation}

\paragraph{Fluctuation--dissipation theorem (FDT).}
A general statement (Callen \& Welton 1951; for magnetism specifically
Brown 1963, \emph{Phys.\ Rev.}\ \textbf{130}, 1677) that connects the
spectral density of thermal fluctuations of a variable to the
dissipation in the system's response. For the stochastic LLG it pins
the noise amplitude to the Gilbert damping $\alpha$: the same friction
that dissipates magnetic energy must also drive the thermal kicks at
temperature $T$, so that the equilibrium distribution is Boltzmann. If
$\sigma$ is set independently of $\alpha$, the system does \emph{not}
thermalise to the Boltzmann distribution; the macrospin Langevin
validation gate (Sec.~\ref{sec:val}) checks precisely this.

Magnetic fluctuations at finite temperature are modeled by augmenting
the deterministic effective field with a Gaussian white-noise field
$\hth(\mb{r}, t)$ that satisfies the fluctuation--dissipation theorem:
%
\begin{equation}\label{eq:sllg}
\frac{d\mb{m}}{dt} = -\gamma' \, \mb{m} \times (\Heff + \hth)
- \gamma' \alpha \, \mb{m} \times \big[\mb{m} \times (\Heff + \hth)\big]
+ \bm{\tau}_{\text{SOT}},
\end{equation}
%
with $\gamma' = \gamma / (1 + \alpha^2)$. The thermal field has zero
mean and is delta-correlated in space and time:
%
\begin{equation}\label{eq:fdt}
\langle h_{\text{th},i}^{(\ell)}(\mb{r}, t)\, h_{\text{th},j}^{(\ell')}(\mb{r}', t') \rangle
= \sigma^2 \, \delta_{ij}\, \delta_{\ell\ell'}\,
\delta(\mb{r} - \mb{r}')\, \delta(t - t'),
\end{equation}
%
with $\ell \in \{\text{top}, \text{bot}\}$ indexing the two SAF
layers. The noise variance is fixed by the FDT in the form
appropriate to the simulator's units (effective field in Tesla,
$\gamma$ in rad/s/T):
%
\begin{equation}\label{eq:sigma}
\boxed{\;
\sigma^2 = \frac{2\, \alpha\, \kB T}{\gamma\, \Ms\, \Vcell},
\qquad
\Vcell = a^2 \, \tCo.
\;}
\end{equation}
%
A short derivation of~\eqref{eq:sigma} is given in
Appendix~\ref{app:fdt}. The two layers of the SAF have independent
thermal baths.

\subsection{Stratonovich interpretation}

The noise in~\eqref{eq:sllg} multiplies $\mb{m}$ through the cross
product, so this is a multiplicative-noise SDE. The constraint
$|\mb{m}|=1$ is preserved only under the \emph{Stratonovich}
interpretation; an It\^{o} interpretation introduces a spurious drift
that violates the norm. We adopt Stratonovich throughout. The
practical consequence is that any midpoint or predictor-corrector
scheme that evaluates the noise term at both the predictor and
corrector stages \emph{with the same noise sample} converges to the
Stratonovich solution, and renormalization at each step keeps the
constraint to machine precision.

\subsection{Drive: SOT plus Joule heating}

The drive is the existing Slonczewski damping-like plus field-like
spin--orbit torque encoded in \texttt{src/simulator/integrator.py}:
%
\begin{equation}
\bm{\tau}_{\text{SOT}}(t) =
\gamma'\, \big[ (H_{DL}(t) + \alpha H_{FL}(t))\, \mb{m} \times \mb{s}
+ (H_{FL}(t) - \alpha H_{DL}(t))\, \mb{s} \big],
\end{equation}
%
with $\mb{s} = (\hat{\mb{p}} \times \mb{m})/|\hat{\mb{p}} \times \mb{m}|$,
$H_{DL}(t) = c_{DL}\, \jcur(t)$, $H_{FL}(t) = c_{FL}\, \jcur(t)$.
The time-varying current $\jcur(t)$ is supplied by
\texttt{p.pulse(t)} (any object implementing \texttt{\_\_call\_\_(t)};
defaults to \texttt{ConstantPulse(p.J\_current)}). The stochastic
stepper threads the substage time \texttt{t} through to
\texttt{llgs\_rhs(m, H, p, t)} so the SOT amplitude tracks the pulse
at every Heun stage. Setting \texttt{p.pulse = ConstantPulse(0.0)}
during the relaxation block forces $\jcur = 0$ without changing
$H_{DL}/H_{FL}$ by hand.

Joule self-heating raises the lattice temperature above the substrate
temperature:
%
\begin{equation}\label{eq:joule}
T(\jcur) = T_{\text{sub}} + \Rth\, \jcur^2,
\end{equation}
%
where $\Rth$ is a stack-dependent thermal resistance. We carry $\Rth$
as a free, user-specified parameter. No literature-anchored value for
this Pt/Co/Ru SAF is assumed; sensitivity to $\Rth$ is part of the
final output. Through~\eqref{eq:joule}, increasing $\jcur$ both
strengthens the drive and raises $T$, so the trade-off in Q4 is
non-trivial.

\subsection{Observables}

The trajectory worker records both layers; downstream analyzers
consume either the top-layer subset (Q1/Q2/Q4) or the cross-layer
diagnostics (Q3a).

\begin{itemize}[leftmargin=2em]
    \item \textbf{Drift velocity} (top \emph{and} bot)
    $\mb{v}_\ell = \langle d\mb{r}_{c,\ell} / dt \rangle$ from
    \texttt{skyrmion\_center\_pbc} unwrapped against periodic
    boundaries (\texttt{unwrap\_trajectory}).
    \item \textbf{Hall angle} (top and bot)
    $\theta_{H,\ell} = \arctan(\langle v_{y,\ell}\rangle / \langle v_{x,\ell}\rangle)$
    fit on the second half of the trajectory.
    \item \textbf{Transverse variance}
    $\sigma_y^2(t) = \langle (y_c(t) - \langle y_c(t)\rangle)^2 \rangle$.
    \item \textbf{Survival probability}
    $P_{\text{surv}}(T, \jcur) = \Pr(|Q(t_{\max})| > q_{\text{th}}
    \text{ for } k_{\text{c}} \text{ consecutive samples})$
    across the ensemble; the K-consecutive guard suppresses false
    positives from finite-$T$ $Q$ noise.
    \item \textbf{Topological charge per layer} $Q_{\text{top}}$,
    $Q_{\text{bot}}$ and the compensation residual
    $Q_{\text{top}} + Q_{\text{bot}}$ (Q3a).
    \item \textbf{Per-layer diameter} $d_{\text{top}}$,
    $d_{\text{bot}}$ from \texttt{skyrmion\_diameter} with the
    appropriate \texttt{core\_polarity} (Q3a).
    \item \textbf{Top-bot lateral offset}
    $\langle |\mb{r}_{c,\text{top}} - \mb{r}_{c,\text{bot}}| \rangle$
    in steady state (Q3a, inter-layer drag).
    \item \textbf{Pair potential} $V(r)$ from short-time forced-pair
    runs at $\jcur = 0$, $T = 300$~K (Q3b).
\end{itemize}

%======================================================================
\section{Numerical scheme}
%======================================================================

\subsection{Stochastic Heun predictor--corrector}

For Stratonovich SDEs with multiplicative noise, the stochastic Heun
scheme is the standard order-1.0 method. For the SAF pair
$(\mb{m}_t, \mb{m}_b)$ at step $n$, per layer:
%
\begin{enumerate}[leftmargin=2em]
    \item Sample a noise increment $\hth$ for each lattice cell with
    cell-wise standard deviation $\sigma / \sqrt{\Delta t}$. The
    \emph{same} sample is reused in both stages below.
    \item Compute $\Heff$ from the existing
    \texttt{effective\_field} (no noise inside the field assembly).
    Form $H_{\text{total}} = \Heff + \hth$ and call
    $f \equiv \texttt{llgs\_rhs}(\mb{m}, H_{\text{total}}, p)$.
    The $\mb{m}\times(\cdot)$ and $\mb{m}\times(\mb{m}\times(\cdot))$
    structure inside \texttt{llgs\_rhs} distributes the noise across
    the precession and damping terms, matching PysLLG's \texttt{Gfunc}
    form $\propto -\mb{m}\times(\hth + \alpha\, \mb{m}\times\hth)$.
    \item Predictor: $\tilde{\mb{m}} = \mb{m} + \Delta t \, f$.
    \emph{Do not} renormalize the predictor.
    \item Corrector (same noise sample, $\Heff$ recomputed at
    $\tilde{\mb{m}}$):
    $\mb{m}^{n+1} = \mb{m} + \tfrac{\Delta t}{2}
    \big[ f(\mb{m}, \Heff(\mb{m}) + \hth) +
          f(\tilde{\mb{m}}, \Heff(\tilde{\mb{m}}) + \hth) \big]$.
    \item Before normalizing, log
    $\max_{i,j} \big|\,\|\mb{m}^{n+1}_{ij}\| - 1\,\big|$ and raise
    \texttt{RuntimeError} if it exceeds a user-supplied tolerance
    \texttt{tol\_norm}. Then renormalize once.
\end{enumerate}
%
Single end-of-step renormalization matches PysLLG and the standard
Stratonovich-Heun derivation; renormalizing the predictor introduces
an $O(\sigma^2 \Delta t)$ bias to the corrector evaluation that the
single-end version avoids. The two SAF layers carry independent
noise draws; within a step, each layer's draw is shared by its own
predictor and corrector.

\subsection{Time step}

We retain the deterministic $\Delta t = 5\times 10^{-14}$~s as a
starting point and verify a posteriori that the equilibrium
$\langle m_z^2 \rangle$ matches the Langevin prediction to within
5\%. If it does not, halve $\Delta t$ and re-run the Langevin
validation. No automatic step-size adaptation is used; the time step
is user-supplied.

\subsection{Reproducibility}

Each trajectory consumes one explicit \texttt{numpy.random.Generator}
constructed from a user-supplied integer seed. Ensembles fold the
trajectory index into the seed (e.g., \texttt{seed\_base + traj\_idx})
so that any single trajectory is byte-reproducible without rerunning
the rest of the ensemble.

%======================================================================
\section{Implementation plan}
%======================================================================

\subsection{Scope and constraints}

\begin{itemize}[leftmargin=2em]
    \item All new code lives in \texttt{src/stochastic\_llgs/}.
    Existing files are \emph{read} and imported but never modified.
    \item Reuse \texttt{effective\_field} from
    \texttt{src/simulator/fields.py}, \texttt{llgs\_rhs} from
    \texttt{src/simulator/integrator.py}, and
    \texttt{effective\_field\_demag\_pair} from
    \texttt{src/phase\_diagram/fields\_demag.py}. Reuse
    \texttt{skyrmion\_center} and \texttt{skyrmion\_diameter} from
    \texttt{src/simulator/analysis.py}, \texttt{topological\_charge}
    from \texttt{src/simulator/main.py}, and \texttt{saf\_skyrmion}
    from \texttt{src/simulator/initial\_conditions.py}.
    \item No \texttt{argparse}. Every runnable module has a
    \texttt{main()} that opens with a labeled User Configuration
    block of snake\_case Python variables.
    \item No default argument values anywhere. Every function
    parameter is provided explicitly by the caller. Missing or
    invalid inputs raise \texttt{RuntimeError} with a precise message
    naming the offending field.
    \item No silent failure paths. No \texttt{try/except} that logs
    and continues. No \texttt{dict.get(key, fallback)}. The
    end-of-step norm-drift check logs
    $\max\big|\,\|\mb{m}\| - 1\,\big|$ and raises
    \texttt{RuntimeError} if it exceeds a user-supplied
    \texttt{tol\_norm}, instead of letting renormalization mask a
    too-coarse $\Delta t$.
    \item Units SI/Tesla throughout.
    \item Backend strategy is \emph{hybrid}: NumPy stochastic Heun
    first (Phase 1--5). A JAX/\texttt{diffrax} port is implemented
    only if Phase 5 wall time exceeds the working budget (Phase 6,
    optional).
\end{itemize}

\subsection{File layout (new code only)}

\begin{lstlisting}
src/stochastic_llgs/
  __init__.py
  parameters_thermal.py     # attach T, R_th, sigma to a p that
                            # the caller has already populated
  thermal_field.py          # sample_thermal_field(rng, shape,
                            #                     sigma, dt)
  integrator_sllg.py        # heun_stochastic_step for the SAF
                            # pair, with optional demag kernels
  joule_heating.py          # T_of_j(j, T_sub, R_th)
  diagnostics.py            # skyrmion_center_pbc,
                            # unwrap_trajectory,
                            # detect_annihilation, hall_angle
  ensemble.py               # ProcessPoolExecutor ensemble driver
  io.py                     # NPZ writers
  validation/
    __init__.py
    macrospin.py
    test_langevin.py
    test_brown_reversal.py
    test_equipartition.py
    test_t0_limit_nodemag.py
    test_t0_limit_demag.py
  production/
    __init__.py
    run_single.py            # trajectory_worker(config)
    scan_tj.py               # Q1, Q2, Q4-driven
    scan_arrhenius.py        # Q4-intrinsic (j=0, tau(T))
    scan_radius.py           # Q3a (inter-layer + size)
    pair_potential.py        # Q3b (V(r), two SAFs)
  jax_backend/               # Phase 6 only; absent until needed
    fields_jax.py
    integrator_diffrax.py

scripts/
  submit_sllg_scan_tj.sh        # SLURM array, 1 cpu/task,
  submit_sllg_scan_arrhenius.sh #     %1000 throttle
  submit_sllg_scan_radius.sh    # 8 cpus/task (scipy.fft demag)
  submit_sllg_pair_potential.sh # 8 cpus/task
  submit_sllg_all.sh            # dispatches all four
  aggregate_sllg.py             # walks per-traj NPZs ->
                                # aggregate.npz per scan
  analyze_sllg_scan_tj.py       # reads aggregate.npz ->
  analyze_sllg_scan_arrhenius.py#   plot + fit + per-scan
  analyze_sllg_scan_radius.py   #   answer NPZ
  analyze_sllg_pair_potential.py
\end{lstlisting}

\subsection{Phase 1 --- Noise generator and macrospin solver}

\begin{enumerate}[leftmargin=2em]
    \item \texttt{parameters\_thermal.attach\_thermal(p, T, R\_th, seed)}
    attaches \texttt{p.T, p.R\_th, p.seed, p.V\_cell, p.k\_B,
    p.sigma\_noise}. All four arguments required; validate finite and
    positive; raise \texttt{RuntimeError} naming any offender.
    \item \texttt{thermal\_field.sample\_thermal\_field(rng, shape,
    sigma, dt)} returns Gaussian noise of shape
    \texttt{(ny, nx, 3)} with std $\sigma/\sqrt{\Delta t}$. The
    \texttt{rng} is supplied by the caller; constructing one
    silently is forbidden.
    \item \texttt{integrator\_sllg.heun\_stochastic\_step(m\_top,
    m\_bot, dt, p, kernels, h\_top, h\_bot, tol\_norm, *, t=0.0)}.
    The thermal field is \emph{pre-sampled} by the caller (via
    \texttt{sample\_thermal\_field}) and threaded in as
    \texttt{h\_top}, \texttt{h\_bot}; the stepper never draws noise
    internally. \texttt{kernels} is \texttt{None} (no demag) or a
    populated dict from \texttt{precompute\_demag\_kernels(p)}.
    \texttt{tol\_norm} is the end-of-step norm-drift cap (raises
    \texttt{RuntimeError} if exceeded). \texttt{t} is the
    keyword-only substage time used by \texttt{llgs\_rhs} to evaluate
    \texttt{p.pulse(t)}; defaults to 0 for constant drives.
\end{enumerate}

Phase 1 gates on the macrospin Langevin check
(Section~\ref{sec:val}).

\subsection{Phase 2 --- Lattice noise and equipartition}

Phase 2 verifies that the noise field has the correct cell-resolved
amplitude. A $16 \times 16$ ferromagnetic patch at $T = 10$~K is
equilibrated, then the mode-resolved energy of low-$k$ magnons is
compared to the Rayleigh--Jeans equipartition $\langle E_k \rangle =
\kB T / 2$. Passing within 20\% on modes with $k < \pi/(4a)$ is the
gate.

\subsection{Phase 3 --- SAF, $T=0$ reproducibility (split)}

The stochastic stepper accepts an optional
\texttt{demag\_kernels} argument. When \texttt{None}, the stepper
calls \texttt{effective\_field} (no demag) for each layer. When a
populated dict from \texttt{precompute\_demag\_kernels(p)} is
supplied, the stepper calls \texttt{effective\_field\_demag\_pair}
from \texttt{src/phase\_diagram/fields\_demag.py}.

\textbf{Phase 3a (no-demag baseline).} Set $T=0$,
\texttt{demag\_kernels=None}, and run the same SOT-driven analysis
as \texttt{src/simulator/analysis.py:run\_analysis}. \textbf{Pass:}
equilibrium diameter, drift velocity, Hall angle, $Q$ each agree
with the deterministic RK4 to within 2\%.

\textbf{Phase 3b (demag pathway).} Set $T=0$ and pass
\texttt{precompute\_demag\_kernels(p)}. The reference is generated
on the fly by driving the existing \texttt{rk4\_step} with
\texttt{effective\_field\_demag\_pair} for the same duration; no
existing code is modified. \textbf{Pass:} 2\% on each observable.

\subsection{Phase 4 --- Diagnostics and Joule heating}

\begin{itemize}[leftmargin=2em]
    \item \texttt{diagnostics.skyrmion\_center\_pbc(m, a,
    core\_polarity)} computes the per-snapshot PBC-aware center by
    the phase of the weighted circular mean
    $\langle e^{i 2\pi x_c / L}\rangle$ over the $m_z$-selected core
    region. \texttt{core\_polarity = +1} for the top layer (core
    $m_z = -1$), $-1$ for the bot layer (core $m_z = +1$).
    \item \texttt{diagnostics.unwrap\_trajectory(cx, cy, L\_x, L\_y)}
    removes $L$-jumps from a wrapped trajectory by detecting
    inter-sample displacements $> L/2$ and adding $\pm L$. Requires
    inter-sample displacement $< L/2$ to be unambiguous.
    \item \texttt{diagnostics.detect\_annihilation(Q\_history,
    q\_threshold, k\_consecutive)} requires both
    \texttt{q\_threshold} and \texttt{k\_consecutive} from the
    caller. Returns the first step index $n$ such that
    $|Q[m]| < $\,\texttt{q\_threshold} holds for all $m \in
    [n - k_c + 1, n]$, or $-1$ if the skyrmion survives. The
    consecutive-step guard suppresses false positives from
    finite-$T$ topological-charge fluctuations.
    \item \texttt{diagnostics.hall\_angle(t, cx, cy, half)} fits a
    linear drift on the last \texttt{half} fraction of the unwrapped
    trajectory and returns $(v_x, v_y, \theta_H)$.
    \item \texttt{joule\_heating.T\_of\_j(j, T\_sub, R\_th)}
    implements~\eqref{eq:joule} as a uniform global temperature
    (one scalar $T$ per run); spatially resolved heating is not
    modeled in this phase.
\end{itemize}

\subsection{Phase 5 --- Production runs (NumPy)}

Shared configuration: lattice $256 \times 256$,
$\Delta t = 5\times 10^{-14}$~s. Joule heating is uniform and
global, $T = T_{\text{sub}} + R_{\text{th}} \jcur^2$. Material
defaults set in \texttt{src/simulator/parameters.py}: $D = 0.85
\times 10^{-3}$~J/m$^2$ (updated from the Pham 2024 Set A value
of $0.62 \times 10^{-3}$), $R_{\text{sk}} = 93.25$~nm, $dw = 27$~nm.
All $T$ values held to $T \le 350$~K in the user-config blocks (see
Sec.~\ref{sec:risks}). $R_{\text{th}} = 0$ by default so
$T_{\text{eff}} = T_{\text{sub}}$; set explicitly if a calibrated
value is available.

\begin{itemize}[leftmargin=2em]
    \item \texttt{run\_single.py}: one $(T_{\text{sub}}, \jcur)$
    trajectory. Demag is a user-set boolean in the User
    Configuration block. Exports
    \texttt{trajectory\_worker(config)} for the scan drivers.

    \item \texttt{scan\_tj.py} (Q1, Q2, Q4-driven): \textbf{demag
    OFF}, $N_{\text{ens}} = 30$.
    Grid $T_{\text{sub}} \in \{10, 25, 50, 100, 150, 200, 250, 300,
    350\}$~K $\times$
    $\jcur \in \{1, 2, 4, 8, 16\} \times 10^{11}$~A/m$^2$.
    Per cell: $\langle v\rangle, \sigma_v, \theta_H,
    \sigma_{\theta_H}, \sigma_y$, and
    $P_{\text{surv}}$ using the \texttt{detect\_annihilation}
    guard. Total $= 9 \times 5 \times 30 = 1350$ trajectories,
    $t_{\max} = 2$~ns each.

    \item \texttt{scan\_arrhenius.py} (Q4-intrinsic, dedicated):
    demag OFF, $\jcur = 0$, $N_{\text{ens}} = 200$. Sweep
    $T \in \{10, 25, 50, 100, 150, 200, 250, 300, 350\}$~K with
    $t_{\max} = 10$~ns. Fit $\log \tau$ vs $1/T$ via the censored
    exponential MLE and report $\Delta E_{\text{barrier}}$ and the
    Brown prefactor $\tau_0$. Total $= 9 \times 200 = 1800$
    trajectories. With the $T \le 350$~K cap, observed annihilation
    rates may be too low to produce a defensible fit; the script
    reports \texttt{nan} for $\Delta E$ in that case.

    \item \texttt{scan\_radius.py} (Q3a, inter-layer + size):
    \textbf{demag ON}, $N_{\text{ens}} = 30$. At $T = 300$~K and
    $\jcur = 4 \times 10^{11}$~A/m$^2$, sweep $(D, H_z)$ over the
    \texttt{dh\_cells} list (defaults inherited from the D=0.62 era;
    retune for D=0.85 if needed). For each cell, alongside top-layer
    $\langle d\rangle, \langle v\rangle, \langle \theta_H\rangle,
    P_{\text{surv}}$, also report the inter-layer Q3a diagnostics
    $\langle |\mb{r}_{c,\text{top}} - \mb{r}_{c,\text{bot}}|\rangle$
    and $\langle Q_{\text{top}} + Q_{\text{bot}}\rangle$ (the
    compensation residual). Total $= 4 \times 30 = 120$ trajectories.

    \item \texttt{pair\_potential.py} (Q3b, V(r)): \textbf{demag
    ON}, $N_{\text{ens}} = 30$. Forced-pair runs at $T = 300$~K,
    $\jcur = 0$, sweeping initial separations
    $r_{\text{init}} \in \{180, 240, 320, 400\}$~nm. Reconstruct
    $V(r)$ from short-time radial force measurements. Total
    $= 4 \times 30 = 120$ trajectories.
\end{itemize}

\paragraph{Why demag-off for the stability scans, demag-on for
radius.} The thin-film shape-anisotropy correction $-\mu_0 \Ms$ is
already folded into the precomputed $C_{\text{anis}}$ in
\texttt{parameters.py:132}, capturing the leading-order
uniform-magnetization demag. The non-uniform residual from
in-plane spin divergence in the domain wall shifts the skyrmion
radius by 5--10\%, $\langle v \rangle$ and $\theta_H$ by $<5\%$,
and does not change the qualitative trend of $P_{\text{surv}}(T,
\jcur)$ or $\tau(T)$. For Q3 the radius itself is the observable,
so the demag pathway is enabled there.

\subsection{Phase 6 --- JAX/diffrax port (only if needed)}

Triggered only if the Phase 5 wall time on the production grid
exceeds the user's working budget. Port \texttt{effective\_field}
and the demag convolution to JAX; use \texttt{diffrax.MultiTerm}
with the \texttt{GeneralShARK} solver and a
\texttt{VirtualBrownianTree} keyed by seed; \texttt{jax.vmap} over
the ensemble axis. Regress against a single NumPy cell before
swapping.

%======================================================================
\section{Validation gates}\label{sec:val}
%======================================================================

Every gate must produce its named NPZ artifact and pass its
quantitative criterion before the next phase begins.

\subsection{Langevin macrospin (Phase 1)}

Single spin in a Zeeman field $B \hat{z}$, no anisotropy, $\alpha=1$.
For each $T \in \{10, 25, 50, 100, 150, 200, 250, 300, 350\}$~K
and several $B$ values, run $\geq 200$ trajectories of length
$\gtrsim 100 / (\alpha \gamma B)$. Compare the time- and
ensemble-averaged $\langle m_z \rangle$ to the Langevin function
$L(\mu B / \kB T)$ with $\mu = \Ms \Vcell$.
\textbf{Pass:} RMS$(\langle m_z\rangle_{\text{sim}} - L)/L < 0.05$.

\subsection{Brown reversal time (Phase 1)}

Uniaxial macrospin with $K = K_{\text{top}}$, no field. For
temperatures yielding $K\Vcell/\kB T \in \{3, 5, 8\}$ (higher
barriers are uncomputably slow), measure the mean first-passage time
across $\geq 100$ trajectories and compare to Brown's high-barrier
formula
%
\begin{equation}
\tau \;\sim\; \frac{(1+\alpha^2)}{\alpha\gamma}\,
\sqrt{\frac{\pi \kB T}{K \Vcell}}\,
\exp\!\Big(\frac{K \Vcell}{\kB T}\Big).
\end{equation}
%
\textbf{Pass:} $|\log_{10}(\tau_{\text{sim}} / \tau_{\text{Brown}})|
< 0.5$ at each barrier.

\subsection{Spin-wave equipartition (Phase 2)}

$16 \times 16$ ferromagnetic patch, single layer (RKKY $= 0$), no
DMI, $K = 0$, weak $H_z$. Equilibrate at $T = 10$~K, then accumulate
$\langle |\delta m_k|^2 \omega_k^2 \rangle$. \textbf{Pass:} match to
$\kB T / 2$ within 20\% for modes with $k < \pi/(4a)$.

\subsection{$T=0$ deterministic limit (Phase 3a, no demag)}

Set $T = 0$, \texttt{demag\_kernels = None}, and run the same
SOT-driven analysis as \texttt{src/simulator/analysis.py:run\_analysis}.
\textbf{Pass:} equilibrium diameter, drift velocity, Hall angle, and
$Q$ each agree with the deterministic RK4 to within 2\%.

\subsection{$T=0$ deterministic limit (Phase 3b, with demag)}

With demag on, the canonical \texttt{saf\_skyrmion} IC at default
$(D, H_z = 0)$ is unstable: during the relaxation block it decays
into a multi-soliton $Q \approx +3$ texture. To probe \emph{stepper
consistency} cleanly (separate from physical stability of the IC)
we apply a small stabilising perpendicular field $H_z = +0.2$~T,
shrink the IC to $R = 40$~nm with $dw = 15$~nm, and shorten the
trajectory to $100$~ps relax $+\,100$~ps drive. The reference is
generated on the fly by driving the existing \texttt{rk4\_step}
with \texttt{effective\_field\_demag\_pair} for the same duration;
no edits to existing code. Both integrators converge to the same
equilibrium morphology in this short window. \textbf{Pass:}
diameter and $Q$ each within 2\%. The per-cell m-field RMS and max
deviation are recorded as diagnostics but not gated: in the
chaotic regime that the IC enters, small $O(\Delta t^2)$
discretisation differences amplify into large local m-field
deviations even though the morphology matches.

%======================================================================
\section{Run-script conventions (no argparse)}
%======================================================================

Every runnable module's \texttt{main()} opens with a labeled User
Configuration block of snake\_case Python variables. The user edits
the file to reconfigure a run. SLURM array indices and the in-process
parallelism are read from environment variables
(\texttt{SLURM\_ARRAY\_TASK\_ID}, \texttt{SWEEP\_NPROC}); a missing
variable raises \texttt{KeyError} rather than silently defaulting.

\begin{lstlisting}[language=python]
def main():
    # =========== User Configuration ===========
    t_sub_list  = [10., 25., 50., 100., 150.,
                   200., 250., 300., 350.]  # K
    j_list      = [1e11, 2e11, 4e11, 8e11, 1.6e12]
    r_th        = 0.0
    t_max       = 2.0e-9
    dt          = 5.0e-14
    seed_base   = 101
    n_ens       = 30
    use_demag   = False
    nx, ny      = 256, 256
    # ========= End User Configuration =========
    ...
\end{lstlisting}

\paragraph{SLURM array dispatch.} Each production script builds a
flat trajectory grid (one entry per (cell, ens\_idx)). At run time
it reads \texttt{SLURM\_ARRAY\_TASK\_ID}; if set, the script
restricts the grid to that single index and runs exactly one
trajectory. Out-of-range indices raise \texttt{RuntimeError}. This
matches the pattern of \texttt{scripts/sweep\_S41\_v\_time.py} and
the other figure-reproduction sweeps. Each task writes its own
NPZ to \texttt{output/stochastic\_llgs/<scan>/}, named by the cell
coordinates and the ensemble index.

\paragraph{Submit scripts.} One \texttt{scripts/submit\_sllg\_<scan>.sh}
per scan, plus a master \texttt{submit\_sllg\_all.sh} that dispatches
all four. Per-task resource layout:
\begin{itemize}[leftmargin=2em, nosep]
    \item \texttt{scan\_tj}, \texttt{scan\_arrhenius}: demag OFF, no
    FFT to thread, \texttt{--cpus-per-task=1} with
    \texttt{--array=0-N\%1000} concurrency cap (N = 1349, 1799).
    \item \texttt{scan\_radius}, \texttt{pair\_potential}: demag ON,
    \texttt{scipy.fft.fft2(workers=-1)} threading in
    \texttt{src/simulator/demag.py}, so
    \texttt{--cpus-per-task=8} and no throttle (N = 119).
\end{itemize}
All four set \texttt{SWEEP\_NPROC=1} (one trajectory per task; the
internal \texttt{mp.Pool} branch is only used for local non-array
runs).

\paragraph{Aggregator and analyzers.} Post-cluster pipeline:
\begin{enumerate}[leftmargin=2em, nosep]
    \item Pull per-trajectory NPZs from the cluster
    (\texttt{rsync ... <cluster>:<scratch>/skyrmion/output/}
    \texttt{stochastic\_llgs/}).
    \item \texttt{python -m scripts.aggregate\_sllg <scan>} walks
    the per-task files and writes
    \texttt{aggregate.npz} (one script, one positional argument
    selecting the scan; no argparse).
    \item \texttt{python -m scripts.analyze\_sllg\_<scan>} reads
    the aggregate, produces the per-question figure and the
    answer NPZ (Pareto front, Arrhenius fit, $v(R)$ +
    inter-layer offset, $V(r)$).
\end{enumerate}

%======================================================================
\section{Verification plan (end-to-end)}
%======================================================================

\begin{enumerate}[leftmargin=2em]
    \item \textbf{Langevin gate.} Run
    \texttt{validation/test\_langevin.py}; confirm RMS criterion.
    \item \textbf{Brown gate.}
    \texttt{validation/test\_brown\_reversal.py};
    confirm log-time criterion.
    \item \textbf{Equipartition gate.}
    \texttt{validation/test\_equipartition.py};
    confirm low-$k$ modes within 20\%.
    \item \textbf{$T=0$ gates.}
    \texttt{validation/test\_t0\_limit\_nodemag.py} and
    \texttt{validation/test\_t0\_limit\_demag.py}; confirm 2\%
    agreement against their respective deterministic references.
    \item \textbf{Single-cell production.}
    \texttt{production/run\_single.py} at $(T_{\text{sub}}, \jcur) =
    (300~\text{K}, 4\times 10^{11}~\text{A/m}^2)$. Inspect the
    trajectory; confirm drift, jitter, and survival.
    \item \textbf{Cost calibration.} From the single-cell wall time,
    estimate total cost of the production grid. If above budget,
    trigger Phase 6.
    \item \textbf{Production scans.} Submit
    \texttt{scripts/submit\_sllg\_all.sh} on the cluster.
    \item \textbf{Aggregate + analyze.}
    \texttt{python -m scripts.aggregate\_sllg <scan>} then
    \texttt{python -m scripts.analyze\_sllg\_<scan>} for each scan.
    Per-figure outputs land in \texttt{output/figures\_sllg/}.
    Q4 is answered jointly from \texttt{scan\_tj} (driven survival
    and Pareto front) and \texttt{scan\_arrhenius} (zero-drive
    barrier $\Delta E$).
\end{enumerate}

%======================================================================
\section{Results so far}\label{sec:results}
%======================================================================

A first cluster sweep has been completed; the numbers below are the
direct readouts from
\texttt{output/stochastic\_llgs/scan\_tj/aggregate.npz} (1350
trajectories, $N_\mathrm{ens}=30$ per cell). Q3 and Q4 are reported in
later sections once the LCC pipeline (Sec.~\ref{sec:risks}) and the
high-T Arrhenius extension are run.

%----------------------------------------------------------------------
\subsection{Q1: same drive at room T sustains the $T=0$ velocity}
%----------------------------------------------------------------------

\textbf{Answer: yes — the same drive works, and slightly better at room~T.}

From \texttt{scan\_tj}, the Pareto-optimal cell at every $T_\mathrm{sub}$
(max $\langle v \rangle$ subject to $P_\mathrm{surv} \ge 0.8$) is
$j = 4 \times 10^{11}$~A/m$^2$ at every $T$ from 10~K to 350~K, with
$\langle v \rangle$ rising monotonically:

\begin{table}[h]
\centering
\begin{tabular}{lcccccc}
\toprule
$T_\mathrm{sub}$ [K] & 10 & 50 & 150 & 250 & 300 & 350 \\
\midrule
$\langle v \rangle$ [m/s] at $j = 4 \times 10^{11}$
   & 372 & 371 & 387 & 404 & 421 & 491 \\
SE [m/s]
   & 0.9 & 1.9 & 12.7 & 15.4 & 21.6 & 27.3 \\
\bottomrule
\end{tabular}
\end{table}

At room temperature ($\approx 300$~K) the same $j$ gives
$\langle v \rangle \approx 421$~m/s versus $\approx 372$~m/s at 10~K — a
\textbf{$\sim 13\%$ increase}, not a decrease. The trend is real (well
above the standard error) and is most plausibly a modest
thermal-Hopkinson-like effect: anisotropy softens with $T$, the
skyrmion grows slightly and de-pins from the lattice discretisation,
and the SOT-driven mobility rises. The drive does \emph{not} have to
be reduced to sustain $j = 4 \times 10^{11}$ operation at room~T.
The speed cap is set by survival, not by temperature: at
$j = 8 \times 10^{11}$ and $j = 1.6 \times 10^{12}$ the survival
probability collapses below 20\,\% at every $T$.

%----------------------------------------------------------------------
\subsection{Q2: directional control at high drive}
%----------------------------------------------------------------------

\textbf{Partial answer.} One observable is reliable; the other is
contaminated by the short fit window.

\paragraph{$\sigma_y$ (transverse jitter): reliable.} Steady-state
$\sigma_y$ from \texttt{scan\_tj} at $j = 4 \times 10^{11}$:

\begin{table}[h]
\centering
\begin{tabular}{lcccccc}
\toprule
$T_\mathrm{sub}$ [K] & 10 & 50 & 150 & 250 & 300 & 350 \\
\midrule
$\sigma_y$ [nm] & 4.0 & 7.4 & 21.8 & 27.3 & 50.0 & 77.5 \\
\bottomrule
\end{tabular}
\end{table}

The transverse jitter grows from $\sim 4$~nm at 10~K to $\sim 78$~nm
at 350~K — substantial, but still well below the 256-nm box
half-width and well below the inter-skyrmion separations of interest
($\sim 200$~nm), so the lateral track is not lost.
\textbf{Directional control is preserved} at $j = 4 \times 10^{11}$
up to 350~K. Above that drive ($j = 8 \times 10^{11}$), $\sigma_y$
blows up to 100--300~nm \emph{and} $P_\mathrm{surv}$ collapses
simultaneously: both indicators agree that the drive has crossed
the stability boundary.

\paragraph{Hall angle: not reliable from this dataset.}
$\langle \theta_H \rangle$ bounces from $-87^\circ$ to $+95^\circ$
across cells (e.g.\ at $T = 300$~K it sits at $-11^\circ$,
$+47^\circ$, $-163^\circ$ across the three lowest $j$ values). A
compensated SAF should have $\theta_H \approx 0^\circ$; the noise
in our measurement dwarfs the physical signal. The root cause is
the Hall-angle fit window: 2~ns of drive sampled at 10-ps cadence
gives only $\sim 10$ post-relax samples, short enough that a
single transverse thermal fluctuation can dominate the linear fit.
\textbf{Q2-Hall-angle is INCOMPLETE.} To close it: extend
\texttt{n\_drive} for a longer fit, compute per-trajectory
$\theta_H$ and report the ensemble median rather than the
$\langle v_x \rangle / \langle v_y \rangle$ mean, or switch to a
circular-statistics estimator that handles the $\pm 180^\circ$
wrap correctly.

%======================================================================
\section{Risks and open items}\label{sec:risks}
%======================================================================

\begin{itemize}[leftmargin=2em]
    \item \textbf{$T \le 350$~K cap.} The stochastic-LLGS framework
    assumes constant $M_s$ within a trajectory, valid only for
    $T \ll T_C$ (Sec.~\ref{sec:notation_fixed_T}). No direct $T_C$
    measurement exists for the Pham 2024 stack in any reference in
    \texttt{refs/}: Pham 2024~\cite{pham2024} does not report it;
    Bandiera 2012~\cite{bandiera2012} measures only PMA at room
    temperature with anneals up to 350~$^\circ$C; Cai
    2012~\cite{cai2012} measures Co$_{44}$Pt$_{56}$ \emph{alloy}
    films, not multilayers, and gets $T_C \approx 446{-}733$~K
    depending on thickness; Watanabe 1994~\cite{watanabe1994}
    measures bulk Co-50at\%Pt alloy ($T_C \approx 1100$~K), again
    not the multilayer; Koyama 2015~\cite{koyama2015} is the only
    Pt/Co/Pt \emph{multilayer} study but limited to
    $t_{\text{Co}} = 0.25{-}0.32$~nm with $T_C \approx 100{-}400$~K
    in that sub-monolayer regime. Linear extrapolation of Koyama's
    2000~K/nm slope breaks down well below the Pham
    $t_{\text{Co}} = 1.58$~nm; $T_C$ at that thickness must
    asymptote toward bulk Co (1388~K). A conservative floor
    $T_C \gtrsim 700$~K and the constant-$M_s$ requirement
    $T < 0.5\,T_C$ together motivate the user-config cap
    $T \le 350$~K applied across every production and validation
    script. The cap is not policed inside \texttt{attach\_thermal};
    the user is responsible for choosing $T$ values in the
    user-config blocks.
    \item $\Rth$ has no literature-anchored value for this Pt/Co/Ru
    SAF stack. Default is $\Rth = 0$ (heating disabled); the user
    sets it explicitly when a calibrated value is available.
    \item $N_{\text{ens}} = 30$ for \texttt{scan\_tj},
    \texttt{scan\_radius}, \texttt{pair\_potential};
    $N_{\text{ens}} = 200$ for \texttt{scan\_arrhenius} where the
    censored exponential MLE benefits most from more events. These
    sit at the lower bound of published practice for direct sLLG
    sampling: Litzius 2017 and Woo 2017 use $N = 1{-}10$ for SOT
    velocity, while Hagemeister 2015~\cite{hagemeister2015} and
    Cort\'es-Ortu\~{n}o 2017~\cite{cortes2017} use $N = 100{-}1000$
    for Arrhenius lifetimes.
    \item $\Delta t = 5\times 10^{-14}$~s is borrowed from the
    deterministic simulator; the Langevin gate provides the
    quantitative check that this $\Delta t$ is small enough at the
    capped temperature. The norm-drift guard
    \texttt{tol\_norm}~$=5\times 10^{-3}$ raises loudly if exceeded.
    \item Production cost on $256 \times 256$: the cluster pipeline
    runs all four scans concurrently as SLURM arrays; the
    long-pole is \texttt{scan\_arrhenius} (1800 tasks, 10-ns
    observation window each). The JAX/diffrax port (Phase 6) is
    only triggered if cluster budget proves insufficient.
\end{itemize}

%======================================================================
\appendix
\section{Derivation of the thermal noise variance}\label{app:fdt}
%======================================================================

Starting from the Landau--Lifshitz equation in SI form with
$\Heff$ in Tesla and $\gamma$ in rad/s/T, the magnetization energy
per cell is $E_{\text{cell}} = -\Ms \Vcell \, \mb{m} \cdot \Heff$ so
that the thermal-equilibrium probability density is
$\propto \exp(-E_{\text{cell}}/\kB T)$. The Fokker--Planck operator
associated with~\eqref{eq:sllg} has stationary solution equal to this
Boltzmann distribution if and only if the noise variance
$\sigma^2$ obeys the Einstein--type relation
%
\begin{equation}
\sigma^2 = \frac{2\, \alpha\, \kB T}{\gamma\, \Ms\, \Vcell},
\end{equation}
%
identical to~\eqref{eq:sigma}. This is the standard
fluctuation--dissipation result for Brown's stochastic LLG
equation~\cite{brown1963,garcia1998}. The factor of $\Vcell = a^2
\tCo$ converts the macrospin formula to a per-cell amplitude
appropriate for the discretized lattice; this is why a finer lattice
at the same $T$ gives a smaller per-cell field but the same
integrated-energy fluctuations.

%======================================================================
\begin{thebibliography}{9}

\bibitem{brown1963}
W.~F.~Brown, ``Thermal fluctuations of a single-domain particle,''
\emph{Physical Review} \textbf{130}, 1677 (1963).

\bibitem{garcia1998}
J.~L.~Garc\'ia-Palacios and F.~J.~L\'azaro,
``Langevin-dynamics study of the dynamical properties of small
magnetic particles,''
\emph{Physical Review B} \textbf{58}, 14937 (1998).

\bibitem{pham2024}
V.~T.~Pham et al., ``Fast current-induced skyrmion motion in
synthetic antiferromagnets,'' \emph{Science} (2024).

\bibitem{litzius2017}
K.~Litzius et al., ``Skyrmion Hall effect revealed by direct
time-resolved X-ray microscopy,'' \emph{Nature Physics} \textbf{13},
170 (2017).

\bibitem{rohart2016}
S.~Rohart, J.~Miltat, A.~Thiaville,
``Path to collapse for an isolated N\'eel skyrmion,''
\emph{Physical Review B} \textbf{93}, 214412 (2016).

\bibitem{kidger2021}
P.~Kidger, ``On Neural Differential Equations,'' PhD thesis,
University of Oxford, 2021 (\emph{diffrax} library).

\bibitem{pysllg}
K.~Sung, \texttt{kswng/PysLLG} (GitHub, accessed 2026): reference
Python+Cython implementation of stochastic LLG. The Heun pattern
adopted in Sec.~3.1 mirrors \texttt{LLG\_runge\_kutta.pyx}.

\bibitem{bandiera2012}
S.~Bandiera et al., ``Enhancement of perpendicular magnetic
anisotropy thanks to Pt insertions in synthetic antiferromagnets,''
\emph{Appl.\ Phys.\ Lett.}\ \textbf{101}, 072410 (2012).

\bibitem{koyama2015}
T.~Koyama et al., ``Dependence of Curie temperature on Pt layer
thickness in Co/Pt system,'' \emph{Appl.\ Phys.\ Lett.}\ \textbf{106},
132409 (2015).

\bibitem{cai2012}
W.~Cai et al., ``Curie temperatures of CoPt ultrathin continuous
films,'' \emph{Appl.\ Phys.\ A} \textbf{107}, 519 (2012).

\bibitem{watanabe1994}
K.~Watanabe, T.~Kaneko, S.~Ohnuma, ``Temperature dependence of
magnetic properties in Co--Pt, Fe--Pt and Cr--Pt permanent magnet
alloys,'' \emph{Mater.\ Trans.\ JIM} \textbf{35}, 136 (1994).

\bibitem{hagemeister2015}
J.~Hagemeister, N.~Romming, K.~von Bergmann, E.~Y.~Vedmedenko,
R.~Wiesendanger, ``Stability of single skyrmionic bits,''
\emph{Nat.\ Comm.}\ \textbf{6}, 8455 (2015).

\bibitem{cortes2017}
D.~Cort\'es-Ortu\~{n}o et al., ``Thermal stability and topological
protection of skyrmions in nanotracks,'' \emph{Sci.\ Rep.}\
\textbf{7}, 4060 (2017).

\end{thebibliography}

\end{document}
